\section{Introducción}
La movilidad urbana representa uno de los principales desafíos de las grandes ciudades modernas. El crecimiento sostenido del parque automotor, combinado con la alta densidad poblacional y comercial de los centros urbanos, genera diariamente problemas vinculados con la congestión vehicular, la circulación ineficiente y la dificultad para acceder a espacios de estacionamiento disponibles. En ciudades como la Ciudad Autónoma de Buenos Aires (CABA), una parte considerable del tránsito diario está asociada a conductores que recorren múltiples cuadras buscando estacionamiento, incrementando no solo los tiempos de traslado, sino también el consumo de combustible, la contaminación ambiental y el nivel de estrés urbano.

Paralelamente, existe una gran cantidad de espacios de estacionamiento privados que permanecen desocupados durante extensos períodos de tiempo. Cocheras residenciales, estacionamientos comerciales y espacios pertenecientes a empresas u organizaciones suelen encontrarse subutilizados debido a la falta de herramientas que permitan conectarlos de manera eficiente con potenciales usuarios. Esta situación evidencia una ineficiencia en la distribución y utilización de recursos urbanos ya existentes.

A partir de esta problemática surge ParkLink, una plataforma digital orientada a facilitar la conexión entre conductores que necesitan estacionamiento y propietarios de espacios disponibles que desean monetizar su capacidad ociosa. El objetivo principal del proyecto es centralizar la oferta de estacionamiento mediante una aplicación accesible y sencilla, permitiendo visualizar espacios disponibles en tiempo real, realizar reservas y gestionar pagos de forma segura.

El proyecto se plantea inicialmente sobre un alcance acotado y escalable, priorizando el desarrollo de un Producto Mínimo Viable (MVP) que permita validar la utilidad y viabilidad del sistema en un entorno real. En esta primera etapa se busca desarrollar una solución funcional basada principalmente en reservas manuales y geolocalización, evitando incorporar tecnologías de automatización más complejas —como sistemas inteligentes de detección y predicción vehicular— que podrían incrementar significativamente los costos y riesgos de implementación.

Asimismo, ParkLink no solo busca brindar beneficios a los conductores, reduciendo tiempos de búsqueda y mejorando la experiencia de estacionamiento, sino también generar nuevas oportunidades económicas para propietarios de espacios subutilizados. De esta manera, el proyecto propone un modelo colaborativo que favorece una utilización más eficiente de la infraestructura urbana existente y contribuye al desarrollo de soluciones de movilidad más sostenibles para las grandes ciudades.

Finalmente, el desarrollo del sistema será llevado a cabo mediante una metodología híbrida basada en prácticas ágiles de Scrum y Kanban, permitiendo organizar el trabajo de manera iterativa, adaptable y orientada a la entrega progresiva de funcionalidades.

La definición del proyecto fue resultado de un proceso de generación de ideas y selección entre múltiples alternativas, cuyo análisis detallado se encuentra desarrollado en el \textbf{Anexo 1}.